\section*{Exercises}
\addcontentsline{toc}{section}{Exercises}
The following exercises collect the main proofs and applications of the Fourier-transform and Laplace-transform identities introduced in this chapter.

\renewcommand{\theenumi}{\thechapter.\arabic{enumi}}
\renewcommand{\labelenumi}{\theenumi.}
\begin{enumerate}
    \item\label{ex:fourier-transform-limit} Starting from the Fourier-series formula on $[-\ell,\ell]$, show that as $\ell\to\infty$ the discrete set of wavenumbers becomes dense and the series turns into the Fourier integral. State the limiting form of the coefficient formula and identify the frequency spacing $\Delta\omega$.

    \item\label{ex:fourier-gaussian} For $f(x)=e^{-\alpha x^2}$ with $\alpha>0$, compute its Fourier transform and show that the transform is again a Gaussian. Interpret the effect of increasing $\alpha$ on the width in position space and frequency space.

    \item\label{ex:fourier-derivative} Prove the derivative rule
    \begin{equation*}
        \mathcal{F}\{f'(x)\}=\mathrm{i}\omega\,\mathcal{F}\{f\}(\omega),
    \end{equation*}
    under suitable decay assumptions, and use it to show that the Fourier transform converts differentiation into multiplication by $\mathrm{i}\omega$.

    \item\label{ex:fourier-convolution} Show that the Fourier transform diagonalizes convolution:
    \begin{equation*}
        \mathcal{F}\{f*g\}(\omega)=\sqrt{2\pi}\,F(\omega)G(\omega).
    \end{equation*}
    Explain why this makes the Fourier transform especially useful in solving linear differential equations and filtering problems.

    \item\label{ex:delta-sifting} Prove the sifting property
    \begin{equation*}
        \int_{-\infty}^{\infty}f(x)\delta(x-a)\,\mathrm{d}x=f(a),
    \end{equation*}
    and show that $\delta(ax)=\delta(x)/|a|$ for $a\neq0$. Interpret this result physically as the limit of a sharply localized pulse of unit area.

    \item\label{ex:laplace-derivative} Let $x(p)=\mathcal{L}\{X(t)\}$. Show that
    \begin{equation*}
        \mathcal{L}\{X'(t)\}=p x(p)-X(0),
    \end{equation*}
    and use this to solve the initial-value problem
    \begin{equation*}
        X''+\omega_0^2X=0,
        \qquad X(0)=X_0,\quad X'(0)=0.
    \end{equation*}

    \item\label{ex:laplace-convolution} Prove the convolution theorem for the Laplace transform:
    \begin{equation*}
        \mathcal{L}\{(f*g)(t)\}=\mathcal{L}\{f\}(p)\mathcal{L}\{g\}(p).
    \end{equation*}
    Give one example where this identity simplifies the evaluation of an inverse Laplace transform.
\end{enumerate}
